\documentclass[12pt]{article}
\usepackage[utf8]{inputenc}

\usepackage{mystyle}

\title{Riemann Surfaces}
\author{Joseph Sullivan}
\date{July 2022}

\begin{document}

\maketitle

\section*{Introduction}

These are some notes on Forster's \textit{Lectures on Riemann Surfaces}.

\section{Definitions and Topological Preliminaries}

\begin{defn}
    An $n$-dimensional \textbf{manifold} (for the purposes of Forster) is a Hausdorff topological space $X$ such that every point $a\in X$ has an open neighborhood which is homeomorphic to an open subset of $\R^n$.
\end{defn}

\begin{defn}
    Let $X$ be a two-dimensional manifold. A \textbf{complex chart} on $X$ is a homeomorphism $\phi: U\to V$ of an open subset $U\subseteq X$ onto an open subset $V\subseteq \C$. Two complex charts $\phi_i: U_i\to V_i, i=1,2$ are said to be \textbf{holomorphically compatible} if the map
    \[
        \phi_2\circ \phi_1^{-1} : \phi_1(U_1\cap U_2) \to \phi_2(U_1\cap U_2)
    \]
    is biholomorphic (holomorphic with holomorphic inverse).
    
    A \textbf{complex atlas} on $X$ is a set $\cA=\{\phi_i:U_i \to V_i, i\in I\}$ of charts which are holomorphically compatible and which cover $X$.
    
    Two complex atlases $\cA, \cA'$ on $X$ are \textbf{analytically equivalent} if every chart of $\cA$ is holomorphically compatible with every chart of $\cA'$.
\end{defn}

\begin{rmk}\;
    \begin{enumerate}[(a)]
        \item We can shrink charts and get holomorphically compatible charts.
        \item Analytic equivalence is an equivalence relation, since the composition of two biholomorphisms is a biholomorphism.
    \end{enumerate}
\end{rmk}

\begin{defn}
    A \textbf{complex structure} $\Sigma$ on a two-dimensional manifold is an equivalence class of analytically equivalent complex atlases $\cA$ on $X$.
\end{defn}

\begin{rmk}
    By taking the union of all complex atlases in an equivalence class, we can get a unique maximum complex atlas $\cA^*$. Given $\cA$ representing $\Sigma$, $\cA^*$ is just the collection of all complex charts biholomorphically compatible with every chart of $\cA$.
\end{rmk}

\begin{defn}
    A \textbf{Riemann surface} is a pair $(X, \Sigma)$ of a connected two-dimensional manifold $X$ and a complex structure $\Sigma$ on $X$.
\end{defn}

\begin{eg}
    The \textbf{Riemann sphere} $\P^1$. As a topological space, it is the one point compactification of $\C$, i.e.,
    \[\P^1 \defeq \C \cup \{\infty\},\]
    with open sets being the usual subsets of $\C$ and subsets $(\C\setminus K) \cup \{\infty\}$ for any compact $K\subseteq \C$.
    
    Set
    \begin{align*}
        U_1 &\defeq \P^1 \setminus \{\infty\} = \C\\
        U_2 \defeq \P^1 \setminus \{0\} = \C^* \cup \{\infty\},
    \end{align*}
    and define $\phi_1: U_1\to \C$ as the identity map and $\phi_2: U_2\to \C$ by
    \[
        \phi_2(z) \defeq \begin{cases}
            \frac{1}{z} & z\in \C^*\\
            0 & z=\infty.
        \end{cases}
    \]
    Then, $\phi_1(U_1\cap U_2)=\phi_2(U_1\cap U_2)=\C^*$, and $\phi_2\circ \phi_1^{-1}$ is just the biholomorphism $z\mapsto \frac{1}{z}$.
\end{eg}

\end{document}
